\section{Parameter Displacement under Stochastic Gradient Descent}
\label{sec:displacement}

Loss spikes and oscillatory updates motivate practical diagnostics that can be
computed from gradients already available during training. We develop these
diagnostics by relating parameter displacement, gradient magnitudes and alignment, 
curvature and sampling noise. 
The proposed quantities characterize constructive and destructive regimes, 
and inform possible interventions without the additional overhead of computing Hessians.

Consider the plain SGD update in Eq.~\eqref{eq:sgd_map_stability}, with constant
learning rate $\eta>0$. As before, $\theta_t\in\mathbb{R}^p$ denotes the
parameters after update $t$, and $g_t$ is the mean mini-batch gradient evaluated
at $\theta_{t-1}$. Summing the updates gives the displacement from initialization,
\begin{equation}
  \Delta_T:=\theta_T-\theta_0=-\eta\sum_{t=1}^T g_t.
  \label{eq:displacement}
\end{equation}
For a model trained from scratch, $\theta_0$ represents the intiialized model that has not
yet learned from the training data. After $T$ updates, the trained parameters
$\theta_T$ encode what the model has learned within that training budget.
The displacement $\Delta_T$ captures the net parameter change through which
this learning is expressed. Its squared norm therefore gives a geometric
measure of the magnitude of learned information in parameter space. 
Expanding the squared displacement separates gradient magnitudes from the interactions 
between different steps 
\begin{equation}
\begin{aligned}
  \|\Delta_T\|^2
  ={}&\eta^2\sum_{t=1}^T\|g_t\|^2\\
     &+2\eta^2\sum_{1\leq u<t\leq T}g_u^\top g_t.
\end{aligned}
\label{eq:magnitude-interaction}
\end{equation}
The terms $\|g_t\|^2$ are nonnegative and determine the squared lengths of
individual SGD steps through $\eta^2\|g_t\|^2$. In full-batch GD, a larger
gradient norm indicates a steeper local loss and, therefore, a larger first-order decrease
along the negative gradient. The cross-step terms
$g_u^\top g_t$ describe how these steps interact -- positive inner products
reinforce the displacement, whereas negative inner products cancel
part of it. Large gradient norms can therefore coexist with cancellation along the training trajectory, and
destructive interference (i.e. persistent $g_u^\top g_t < 0$) can limit displacement growth.

To isolate directional alignment from gradient magnitudes, define the cosine
similarity between any two nonzero gradients and its consecutive-step form
\begin{equation}
\begin{aligned}
  \rho_{u,t}&:=\frac{g_u^\top g_t}{\|g_u\|\|g_t\|},
       &&1\leq u<t\leq T.
\end{aligned}
\label{eq:rho}
\end{equation}
In this setup, the alignment between consecutive gradients is
$\rho_t:=\rho_{t-1,t}$, $t=2,\ldots,T$.
Substituting into Eq.~\eqref{eq:magnitude-interaction}
expresses the squared displacement entirely through gradient norms and
directional alignment
\begin{equation}
\begin{aligned}
\|\Delta_T\|^2={}&\eta^2\sum_{t=1}^T\|g_t\|^2\\
 &+2\eta^2\sum_{1\leq u<t\leq T}
    \|g_u\|\|g_t\|\rho_{u,t}.
\end{aligned}
\label{eq:magnitude-alignment}
\end{equation}
For fixed gradient norms, positive $\rho_{u,t}$ increases the squared net
displacement, negative $\rho_{u,t}$ contributes cancellation, and orthogonal
gradients have no pairwise contribution. In practice, 
$\rho_t$ provides an online diagnostic of the
interaction between successive updates: persistent positive or negative
values characterize constructive or destructive regimes.

\subsection{Directional Curvature and Consecutive Alignment}

For mini-batch SGD, we explicitly incorporate stochastic gradient noise in
Eq.~\eqref{eq:sgd_noise_decomposition},
$g_t=\nabla\mathcal{L}(\theta_{t-1})+\xi_t$.
Suppose the full objective $\mathcal{L}$ is twice continuously differentiable
along the preceding step, and write
$H_{t-1}:=\nabla^2\mathcal{L}(\theta_{t-2})$.
Expanding the full gradient around $\theta_{t-2}$ and substituting the preceding
SGD update yields
\begin{equation}
  g_t=(I-\eta H_{t-1})g_{t-1}+(\xi_t-\xi_{t-1})+r_t,
  \label{eq:gradient-evolution}
\end{equation}
where $r_t$ is the Taylor remainder. Neglecting $r_t$ recovers the first-order
approximation. Taking the inner product with $g_{t-1}$ then gives
\begin{equation}
\begin{aligned}
\|g_{t-1}\|\|g_t\|\rho_t
 &\approx\underbrace{\|g_{t-1}\|^2}_{\text{Gradient magnitude}}\\
 &\quad-\underbrace{\eta g_{t-1}^{\top}H_{t-1}g_{t-1}}_{\text{Curvature contribution}}\\
 &\quad+\underbrace{g_{t-1}^{\top}(\xi_t-\xi_{t-1})}_{\text{Noise contribution}}.
\end{aligned}
\label{eq:rho-expanded}
\end{equation}
The right-hand side separates the preceding gradient magnitude,
the curvature contribution, and the gradient noise contribution.
Positive curvature along the preceding gradient direction reduces
reduces its alignment with the next gradient, while
the noise difference can reinforce or oppose this effect.
To characterize these contributions, we define the directional
curvature and normalized noise contribution
\begin{equation}
  \kappa_t:=\frac{g_{t-1}^\top H_{t-1}g_{t-1}}{\|g_{t-1}\|^2},
  \label{eq:kappa}
\end{equation}
\begin{equation}
  \varepsilon_t:=\frac{g_{t-1}^\top(\xi_t-\xi_{t-1})}
                            {\|g_{t-1}\|^2}.
  \label{eq:epsilon}
\end{equation}
Here, $\kappa_t$ measures the curvature projected at the gradient
direction, and $\varepsilon_t$ describes the mini-batch noise alignment with the same gradient direction.
Substituting these into Eq.~\eqref{eq:rho-expanded} and
dividing by $\|g_{t-1}\|\|g_t\|$ gives, after neglecting the Taylor remainder,
the first-order relation
\begin{equation}
  \rho_t\approx\frac{\|g_{t-1}\|}{\|g_t\|}
                   (1-\eta\kappa_t+\varepsilon_t).
  \label{eq:rho-normalized}
\end{equation}
Rearranging Eq.~\eqref{eq:rho-normalized} gives
\begin{equation}
  \eta\widehat\kappa_t
  :=1-\rho_t\frac{\|g_t\|}{\|g_{t-1}\|}
  \approx \eta\kappa_t-\varepsilon_t.
  \label{eq:kappa-estimate}
\end{equation}
The quantity $\widehat\kappa_t$ is computable from consecutive
mini-batch gradients without evaluating the Hessian. However,
it combines directional curvature with stochastic noise and
therefore does not isolate $\kappa_t$.
Equation~\eqref{eq:rho-normalized} also characterizes the sign
of consecutive gradient alignment. Neglecting the Taylor remainder,
and in the absence of noise ($\varepsilon_t=0$), consecutive
nonzero gradients are positively aligned when $\eta\kappa_t<1$,
orthogonal when $\eta\kappa_t=1$, and negatively aligned when
$\eta\kappa_t>1$.
Differently from Eq.~\ref{eq:quad_stability_condition},
here we focus on the caracterization of constructive-destructive regimes.


\subsection{Noise Contribution under Mini-Batch Sampling}
\label{subsec:sampling-noise}

Although mini-batch gradients are unbiased estimators of the full
gradient, the noise contribution to consecutive gradient alignment
need not vanish in expectation. Under fresh, independent mini-batch
sampling across updates, the current noise $\xi_t$ has zero
expected inner product with the preceding gradient $g_{t-1}$.
However, $g_{t-1}$ already contains its own noise $\xi_{t-1}$.
Consequently, assuming finite second moments,
\begin{equation}
  \mathbb{E}[g_{t-1}^\top(\xi_t-\xi_{t-1})]
  =-\mathbb{E}[\|\xi_{t-1}\|^2]\leq0.
  \label{eq:noise-expectation}
\end{equation}
Thus, even unbiased sampling noise contributes a systematic
nonpositive term to the expected unnormalized interaction
between consecutive gradients.

This contribution is directly related to mini-batch size.
Consider a mini-batch of $b$ examples sampled uniformly
without replacement from $N>1$ training examples.
Using the per-sample gradient covariance $\Sigma(\theta)$
defined in Eq.~\eqref{eq:grad_cov_def}, the sampling
variance satisfies
\begin{equation}
  \mathbb{E}[\|\xi_{t-1}\|^2]
  =\frac{N-b}{b(N-1)}
    \mathbb{E}[\operatorname{tr}\Sigma(\theta_{t-2})].
  \label{eq:finite-population-noise}
\end{equation}
At fixed parameters, smaller batches produce larger sampling
variance and, consequently, a larger negative expected noise
contribution. For full-batch gradients ($b=N$), this contribution
vanishes.

Importantly, this does not imply that the total expected inner
product $g_{t-1}^\top g_t$ necessarily becomes more negative
with smaller batches. The same sampling variance increases
the expected squared norm $\|g_{t-1}\|^2$, and these two
contributions cancel in Eq.~\eqref{eq:rho-expanded}.
Nevertheless, increasing noise can reduce directional coherence
by increasing gradient magnitudes relative to their common
full-gradient component.

To illustrate this effect, suppose the full gradient changes
little between consecutive updates, so that both gradients
behave approximately as independent noisy estimates at the
same parameters. When their norms are sufficiently concentrated,
the expected cosine similarity can be approximated by
\begin{equation}
  \mathbb{E}[\rho_t]\approx
  \frac{\|\nabla\mathcal{L}(\theta)\|^2}
  {\|\nabla\mathcal{L}(\theta)\|^2+
  \frac{N-b}{b(N-1)}\operatorname{tr}\Sigma(\theta)}.
  \label{eq:noise-cosine-approx}
\end{equation}
This approximation suggests that, as sampling noise dominates
the full-gradient signal, consecutive gradients become less
directionally coherent, with their cosine similarity approaching
zero. The resulting trajectory becomes increasingly driven by
stochastic fluctuations rather than a consistent descent direction.
However, this behavior is not universal: curvature couples
successive SGD updates and can induce systematic negative
alignment even in noise-dominated regimes. Therefore, batch
size influences observed alignment through both sampling noise
and its interaction with local curvature.


\subsection{Normalized Displacement and Cumulative Interference}
\label{subsec:normalized-displacement}

While the consecutive cosine similarity $\rho_t$ measures local lag-1
directional alignment, it does not capture interactions between
more distant updates (lag$\leq2$). To measure their cumulative effect,
we define the normalized displacement
\begin{equation}
  R_T:=
  \frac{\|\Delta_T\|}
       {\eta\sum_{t=1}^T\|g_t\|}
  =\frac{\|\sum_{t=1}^T g_t\|}
        {\sum_{t=1}^T\|g_t\|},
  \qquad 0\leq R_T\leq1,
  \label{eq:normalized-displacement}
\end{equation}
For constant-step plain
SGD, the denominator represents the total distance traveled
in parameter space, while the numerator measures the net
displacement from initialization. Thus, $R_T$ quantifies the
fraction of accumulated step length retained as net displacement.

Using Eq.~\eqref{eq:magnitude-alignment}, this ratio satisfies
\begin{equation}
\begin{aligned}
R_T^2
=\frac{
\sum_{t=1}^T\|g_t\|^2
+2\sum_{1\leq u<t\leq T}
\|g_u\|\|g_t\|\rho_{u,t}
}{
\left(\sum_{t=1}^T\|g_t\|\right)^2
}.
\end{aligned}
\label{eq:cumulative-interference}
\end{equation}
Unlike $\rho_t$, which considers only adjacent gradients,
$R_T$ incorporates interactions across all temporal lags.
Values near one indicate that successive updates reinforce
a consistent net direction, while smaller values reveal
greater cumulative cancellation.

\subsection{Connection with Loss Spikes and Convergence}
\label{subsec:loss-connection}

Constructive and destructive regimes have been characterized on parameter
motion. To connect gradient alignment with optimization progress, consider
the loss change along the preceding SGD update. A second-order
approximation gives
\begin{equation}
\begin{aligned}
&\mathcal{L}(\theta_{t-1})-\mathcal{L}(\theta_{t-2})\\
&\quad\approx
-\frac{\eta}{2}\|g_{t-1}\|^2
\left(1+\rho_t\frac{\|g_t\|}{\|g_{t-1}\|}\right)\\
&\qquad+\frac{\eta}{2}
g_{t-1}^\top(\xi_{t-1}+\xi_t).
\end{aligned}
\label{eq:loss-curvature}
\end{equation}
The complete derivation is provided at Appendix \ref{appendix:loss-curvature}. 
For full-batch gradient descent, the noise contribution vanishes.
From Eq.~\eqref{eq:rho-normalized}, $\rho_t\|g_t\|/\|g_{t-1}\|\approx 1-\eta\kappa_t$.
Consequently, the local quadratic model predicts a loss decrease
when $\eta\kappa_t<2$, whereas consecutive gradients become
negatively aligned when $\eta\kappa_t>1$. Thus, the interval
$1<\eta\kappa_t<2$ describes destructive interference without
loss increase. Negative alignment alone is therefore insufficient
to characterize a loss spike, once it depends too on the relative gradient magnitudes.

Under mini-batch SGD, the noise contribution can overcome the
predicted decrease, producing a loss spike even when
$\eta\kappa_t<2$. Such an increase may be transient if subsequent
updates recover the preceding loss level. 
